\documentclass[10pt,a4paper]{amsart}
\usepackage[margin=19mm]{geometry}
\usepackage{amsmath,amssymb,mathtools}
\usepackage[scr=rsfs,cal=euler]{mathalfa}
\usepackage{xfrac}
\usepackage[unicode,hidelinks,bookmarks=false]{hyperref}
\newcommand{\ie}{i.e.\ }
\newcommand{\cf}{cf.\ }
\newcommand{\resp}{respectively\ }
\newcommand{\Subsection}{\subsection}
\providecommand{\nobreakdash}{\nobreak\hbox{-}}
\setlength{\emergencystretch}{2em}
\multlinegap0pt
\usepackage{bm}
\usepackage{xcolor}
\usepackage{mathtools}
\usepackage{amssymb}
\usepackage{csquotes}
\usepackage{tikz}
\newtheorem{assumption}{Assumption}
\newtheorem{theorem}{Theorem}
\newcommand{\GI}{{\Gamma_{\mathrm I}}}
\newcommand{\Grad}[1]{\nabla #1}
\newcommand{\LI}{L_{\mathrm I}}
\newcommand{\NormDer}{{\partial_{\bm{n}}}}
\newcommand{\OT}{{\Omega_T}}
\newcommand{\OZ}{{\Omega_0}}
\newcommand{\SI}{{\Sigma_{\mathrm I}}}
\newcommand{\TimeDer}{{\partial^{}_t}}
\newcommand{\conj}[1]{\overline{#1}}
\newcommand{\dS}{{\,\mathrm{d}S}}
\newcommand{\deltaI}{{\delta_{\mathrm I}}}
\newcommand{\dt}{{\,\mathrm{d}t}}
\newcommand{\dx}{{\,\mathrm{d}\vect x}}
\newcommand{\norm}[1]{\|#1\|}
\newcommand{\oW}{{\operator{W}}}
\newcommand{\operator}[1]{\mathcal{#1}}
\newcommand{\vect}[1]{\bm{#1}}
\newcommand{\vx}{{\vect x}}
\title{A space--time continuous and coercive formulation for the wave equation}
\author{Paolo Bignardi, Andrea Moiola}
\date{Source: arXiv:2312.07268v2 (CC BY 4.0)}
\begin{document}
\maketitle

\noindent\emph{Source and licence.} A space--time continuous and coercive formulation for the wave equation. Version arXiv:2312.07268v2 (CC BY 4.0), distributed by arXiv under the Creative Commons Attribution 4.0 International licence, http://creativecommons.org/licenses/by/4.0/, as recorded in the arXiv licence metadata for this exact version and shown on the abstract page https://arxiv.org/abs/2312.07268v2. Full licence notice: \texttt{licenses/arxiv-2312.07268-CC-BY-4.0.txt}.\par\smallskip

\noindent\textit{Editorial scope.} Counted unit: the article's theorem thm:coerc\_b (source0.tex lines 1055--1072). The article's complete original proof of it is delivered with it (source0.tex lines 1073--1185, one \texttt{proof} environment of 113 lines). No mathematics is written, repaired or re-derived: the statement and the proof are the article's own lines, in the article's own notation, and no retained line is substituted or re-flowed.  Retained uncounted as the source-local context the proof needs: lines 664--668; lines 873--893; lines 931--939; lines 1035--1037.  Nothing was omitted from any retained span, so the package carries no omission record.  No retained line contains a citation, so the wrapper carries no bibliography and no bibliography entry.  No retained line contains a figure environment, an image-inclusion command or drawing code, so no image is delivered and the package declares no asset dependency.  Editorial formatting only, in addition: an AMS \texttt{amsart} preamble that restates the article's own declarations and macros used by the retained lines, namely the environments \texttt{assumption}, \texttt{theorem} on separate counters, together with the standard packages those lines need.  The retained environments print in this document as Assumption 1, Theorem 1 in order of appearance, whereas the article prints the same items with the numbering of its own full document.  The complete native source of the article as distributed in the arXiv e-print for this exact version is bundled unchanged as \texttt{sources/151374/source0.tex}.
\begin{align}\begin{aligned}
		\Re b(v,v) & = \ell(v,v) + \Re \int_D \operator{M}v \conj{\operator{L}v} \dx - \Re r(v,v)               \\
		           & = \ell(v, v) + \frac{1}{2} \Re r(v,v) + \frac{1}{2} \Re s(v,v) - \Re r(v,v)                \\
		           & = \ell(v, v) - \frac{1}{2}\Re r(v,v) + \frac{1}{2} \Re s(v,v) \qquad\qquad \forall v\in V.
	\end{aligned}\label{eq:abstract_coercivity_condition}\end{align}

\begin{align}	\label{eq:r_def}
	 & \begin{aligned}
		   r(u,v) := - & \int_Q \left[
			   (\beta + \xi d)\TimeDer{u}\TimeDer{v} + c^2  (\beta + 2 \xi - d\xi)\Grad u \cdot \Grad v
		   \right] \dx \dt                                                                                              \\
+           & \int_{\Omega_T} \big[-\xi \vect x \cdot (  \TimeDer{u} \,\Grad v+ \Grad u\, \TimeDer{v} )
+\beta (T-T^*)(\TimeDer{u} \TimeDer{v} + c^2 \Grad u \cdot \Grad v)\big] \dx                              \\
-           & \int_{\SI} \big[c^2 \operator{M} u\,\NormDer{v} -\frac{c}{\theta} \TimeDer{u}\, \operator{M} v
		   + \xi \,\vect x \cdot \vect{n}\,(-\TimeDer{u} \TimeDer{v}+  c^2\,\Grad u \cdot \Grad v)\big] \dS\dt,   \\
	   \end{aligned}        \\
\label{eq:s_def}
	 & \begin{aligned}
		   s(u,v) := & \int_{\Omega_0} \big[ \xi \vect x \cdot \big( \TimeDer{u} \,\Grad v+ \Grad u \,\TimeDer{v}\big)
		   +\beta T^* \left(\TimeDer{u} \TimeDer{v} + c^2 \Grad u \cdot \Grad v\right)\big] \dx                        \\
& -  \int_{\SI} c^2 \left[(\theta c)^{-1} \TimeDer{u} + \NormDer{u}\right]\operator{M} v  \dS \dt,
	   \end{aligned} \\
	\label{eq:ell_def}
	 & \begin{aligned}
		   \ell(u,v) := \; & A_QT^2\int_Q \oW u \oW v \dx \dt + A_\OZ T^{-1}\int_{\OZ} u v \dx,
	   \end{aligned}
\end{align}

\begin{equation}
	\label{eq:norm_def}
	\begin{split}
		\|v\|_V^2 :=& \qquad\; \|\TimeDer{v}\|_{Q}^2 + \quad c^2\|\Grad v\|_Q^2 + T^2 \|\operator{W} v\|_Q^2 \\
		& + T \|\TimeDer{v}\|_{\Omega_T}^2 + c^2 T \|\Grad v\|_{\Omega_T}^2 \\
		& + T \|\TimeDer{v}\|_{\Omega_0}^2 \:+ c^2 T \|\Grad v\|_{\Omega_0}^2 + T^{-1}\|v\|_{\OZ}^2 \\
		& + \LI \|\TimeDer{v}\|_{\SI}^2 + c^2 \LI \|\Grad v\|_{\SI}^2 .
	\end{split}
\end{equation}

\begin{assumption}[Star-shaped $\Omega$]\label{ass:StarS}
	There exists $\deltaI > 0 $ such that $\vect{x}\cdot\vect{n} \geq \deltaI \LI$ for almost every $\vect x\in \GI$.
\end{assumption}

\begin{theorem}[Coercivity of $b$]\label{thm:coerc_b}
	Under Assumption~\ref{ass:StarS} on $\Omega$, if $\xi > 0$, $A_Q, A_\OZ > 0$ and
	\begin{equation}
		\label{eq:coefficient_coercivity_cond}
		\beta \geq \max
		\begin{cases}
			\xi(d-1),                                               \\
			\xi\frac{1}{\nu - 1}\left(\frac{\LI}{cT}+ 1\right), \\
			\xi\frac{1}{\nu - 1}\frac{\LI}{cT}\left(\theta + \frac1{\deltaI \theta}\right),
		\end{cases}
	\end{equation}
	then for all $u \in W$
	\begin{equation} \label{eq:coercivity_condition}
		b(u,u) \geq {\alpha_b} \|u\|_V^2
		\qquad \text{with}\qquad
		\alpha_b := \min \left\{\frac{\xi\deltaI}{4}, \; A_Q, \; A_\OZ\right\}.
	\end{equation}
\end{theorem}

\begin{proof}
	Let $u \in W$.
	We recall from \eqref{eq:abstract_coercivity_condition} that, in order to prove the coercivity of $b(\cdot,\cdot)$, it is sufficient to show a lower bound for $-r(u,u) + s(u,u) + 2\ell(u,u)$. For simplicity we can find a lower bound for $-r(u,u) + s(u,u)$ and then improve on it by including the terms in $\ell(u,u)$.
	From \eqref{eq:r_def}--\eqref{eq:s_def}, simple calculations lead to
	\begin{align}\nonumber
		-r(u,u) + s(u,u) & =
		\int_{Q} \big[
			(\beta + \xi d) (\TimeDer{u})^2
			+ c^2 ( \beta + 2\xi - d\xi)|\Grad u|^2
		\big] \dx \dt                                           \\ \nonumber
		                 & + \int_{\OT} \big[
			2 \xi (\vect{x} \cdot \Grad u) \TimeDer{u}
			- \beta (T-T^*)\left(
			(\TimeDer{u})^2
			+ c^2 |\Grad u|^2
		\right) \big] \dx                                       \\
		                 & + \int_{\OZ}
		\big[\beta T^* \left(
			(\TimeDer{u})^2
			+ c^2 |\Grad u|^2\right)
			+ 2 \xi (\vect{x} \cdot \Grad u) \TimeDer{u} \big]\dx
		\label{eq:coercivity_step1}                             \\
		                 & + \int_{{\SI}} \left[
			2 \frac c\theta\xi (\vect{x} \cdot \Grad u) \TimeDer{u}
+ c^2 \xi (\vect{x} \cdot \vect{n}) |\Grad u|^2
			- \left(\xi (\vect{x} \cdot \vect{n})
			+ 2\frac{c}{\theta}{\beta} (t-T^*)\right) (\TimeDer{u})^2
			\right] \dS \dt
		\nonumber
		\\
		                 & =: I_Q+I_\OT+I_\OZ+I_{\SI}.\nonumber
	\end{align}
	The mixed terms containing $\Grad u \TimeDer{u}$ can be lower-bounded by negative terms using the H\"older inequality, and then the products of the norms of $\Grad u$ and $\TimeDer{u}$ (on $\OZ$, $\OT$ and ${\SI}$) can be controlled by the weighted Young inequality:
\begin{align*}
		\int_{D}2  (\vect{x} \cdot \Grad u) \TimeDer{u} \ge - 2 \LI \norm{\Grad u}_D \norm{\TimeDer{u}}_D
		 & \geq -\epsilon_D \LI \norm{\Grad u}^2_D -\frac{1}{\epsilon_D} \LI \norm{\TimeDer{u}}^2_D,
		\qquad D\in\{\Omega_0,\Omega_T,{\SI}\},\; \epsilon_D>0.
	\end{align*}
Here we also used that $\LI = \max\{|\vx\|: x\in\Omega\}$.
	Combining these inequalities into \eqref{eq:coercivity_step1} it is possible to determine conditions for the coefficients $\beta$, $\xi$, $T^*$.
	We tackle the different integrals independently.

	\textit{Integrals over $\OT$ and over $\OZ$.}
	Recalling that $T^* = \nu T$, the integrals over $\OT$ and $\OZ$ are bounded from below as follows:
	\begin{align*}
		 & I_\OT\geq \left(
		\beta (\nu-1) - \frac{\xi}{\epsilon_{\Omega_T}}\frac{\LI}{T}
		\right) T \norm{\TimeDer{u}}^2_{\OT}
		+ \left(\beta (\nu-1) - \epsilon_{\Omega_T} \xi \frac{\LI}{c^2 T}
		\right) c^2 T \norm{\Grad u}^2_{\OT},
		\\
		 & I_\OZ\geq \left(
		\beta \nu - \frac{\xi}{\epsilon_{\Omega_0}} \frac{\LI}{T}
		\right) T \norm{\TimeDer{u}}^2_{\OZ}
		+ \left(
		\beta \nu - \epsilon_{\Omega_0} \xi \frac{\LI}{c^2 T}
		\right) c^2 T \norm{\Grad u}^2_{\OZ}.
	\end{align*}
	Taking $\epsilon_{\Omega_T} = \epsilon_{\Omega_0} = c$, the coefficients in the two brackets of each inequality are dimensionless and coincide.
From the assumption \eqref{eq:coefficient_coercivity_cond} on $\beta$,
\begin{equation*}
		\beta \nu - \xi \frac{\LI}{cT} \geq \beta (\nu - 1) - \xi \frac{\LI}{cT} \geq \xi > 0,
	\end{equation*}
	thus
	\begin{equation}\label{eq:coerc:OTOZ}
		I_\OT\geq \xi T \norm{\TimeDer{u}}^2_\OT+\xi c^2 T \norm{\Grad u}^2_\OT,\qquad
		I_\OZ\geq \xi T \norm{\TimeDer{u}}^2_\OZ+\xi c^2 T \norm{\Grad u}^2_\OZ.
	\end{equation}


	\textit{Integral over ${\SI}$.}
	Using the same strategy as for the other integrals and the fact that $\vect{x} \cdot \vect{n} \geq \deltaI \LI$ for all $\vect{x} \in \partial {\Omega}$ by Assumption~\ref{ass:StarS}, it holds that
	\begin{align*}
		I_{\SI} \geq \left(
		\xi \deltaI - \frac{\epsilon_\SI \xi}{\theta} \frac{1}{c}
		\right) c^2 \LI \norm{\Grad u}^2_{{\SI}}
		+ \left(
		\frac2\theta \beta (\nu - 1) \frac{cT}{\LI} - \xi \Big(1 + \frac{c}{\theta \epsilon_{\SI}}\Big)
		\right) \LI \norm{\TimeDer{u}}_{{\SI}}^2.
	\end{align*}
	We take $\epsilon_{\SI} = \frac{\theta \deltaI c}{2}$, so that $ \xi \deltaI - \frac{\epsilon_{\SI}}{\theta c} \xi = \xi\deltaI/2 > 0$. From the the last condition \eqref{eq:coefficient_coercivity_cond} on $\beta$,
	\begin{equation*}
			\frac2\theta  \beta (\nu - 1) \frac{cT}{\LI} - \xi \left(1 + \frac{2}{\deltaI \theta^2} \right)
\geq \xi
	\end{equation*}
	thus
	\begin{equation}\label{eq:coerc:S}
		I_{\SI}\ge \frac{\xi\deltaI}2 c^2 \LI \norm{\Grad u}^2_{{\SI}}  + \xi \LI \norm{\TimeDer{u}}_{{\SI}}^2.
	\end{equation}

	\textit{Integral over ${Q}$.}
	Using again the assumption \eqref{eq:coefficient_coercivity_cond},
	\begin{equation}\label{eq:coerc:Q}
		I_Q \ge \xi\norm{\TimeDer{u}}^2_{{Q}} +\xi c^2\norm{\Grad u}^2_{{Q}}.
	\end{equation}


	\textit{Coercivity constant.}
To obtain the coercivity bound, we recall \eqref{eq:abstract_coercivity_condition}, sum the bounds \eqref{eq:coerc:OTOZ}, \eqref{eq:coerc:S}, and \eqref{eq:coerc:Q}, use that $\deltaI\le1$ and the definition \eqref{eq:norm_def} of the norm $\|\cdot\|_V$:
\begin{align*}
		b(u,u)
		 & = \ell(u,u) -\frac12 r(u,u)+\frac12s(u,u)                                                                      \\
		 & =A_Q T^2\norm{\oW u}_Q^2 + A_\OZ\frac1T\|u\|_\OZ^2+\frac12( I_Q+I_\OT+I_\OZ+I_\SI)                             \\
		 & \ge A_Q T^2\norm{\oW u}_Q^2 + A_\OZ\frac1T\|u\|_\OZ^2+\frac\xi2\Big(
		T \norm{\TimeDer{u}}_{\OT}^2
		+ c^2 T \norm{\Grad u}_{\OT}^2
		+ T \norm{\TimeDer{u}}_{\OZ}^2
		+ c^2 T \norm{\Grad u}_{\OZ}^2                                                                                    \\
		 & \hspace{50mm} +  \LI \norm{\TimeDer{u}}_{{\SI}}^2 + \frac{\deltaI}{2} c^2 \LI \norm{\Grad u}_{{\SI}}^2
		+ \norm{\TimeDer{u}}_{{Q}}^2 + c^2 \norm{\Grad u}_{{Q}}^2\bigg)                                                   \\
		 & \ge \min\Big\{\frac{\xi\deltaI}4, A_Q, A_\OZ\Big\} \norm{u}_V^2.
	\end{align*}
\end{proof}

\end{document}
